\chapter{上半带的公式，先看它们在节省哪种位置}\label{ch:upper}
\section{这一讲还要证明什么}
上一讲说明，合格的计数表会自动接出剪切和同余输送。本讲只负责证明：对每个允许的 $p,s$，真的有这样的表。我们仍记 $r=p+s$，$1\le s<p$。

不必一开始记住所有公式。先按新值 $D$ 还剩多少份来理解三种选择。新载体保留 $p+s$ 份，因此可用于输出的 $D$ 只有 $p-s$ 份。
\begin{enumerate}
\item $s$ 很大、接近 $p$ 时，$D$ 剩得少，直接设计不使用 $D$ 的输出。
\item $s$ 在中段时，让旧三类取用更均衡，用一个整数 $e$ 在两种输出之间调整。
\item $s$ 很小时，用另外两级相邻参数保证根系数的最大公因数足够小。
\end{enumerate}
三类设计的算术收尾相同，不同的是剩余位置怎么安排。
\InsertProof{prop:upper}

\section{一个可行整数不够，需要相邻两个}
典型输出写成 $o_e=(\alpha-e,p-3\alpha-e,2\alpha-s+2e,s)$。当 $e$ 增加一，恰好是一份 $A$、一份 $B$ 换成两份 $C$。这就是上一讲产生剪切的相邻输出。

所有库存限制会变成 $L\le e\le U$。我们必须证明 $L+1\le U$，因为一个可行点只给一个反射，两点才给反射之间的比较。

例如 $p=83,s=10$，取 $\alpha=14$。来源 $(28,14,41)$，剩余库存为 $(27,55,11,73)$。可选 $e=-9,-8$，得到输出
\[(23,50,0,10),qquad(22,49,2,10).
\]
两者都非负、和为83、没有超库存。这里 $e$ 为负完全可以：$e$ 是参数，不是直接取用份数；真正不能为负的是输出的四项。

第二个相邻 $\alpha=15$ 同样给两张表。两个 $\alpha$ 互素；根系数 $H=3n\alpha-2p^2$ 分别为 $-2900,-2123$，最大公因数为一。这就满足上一讲的资源条件。
\InsertProof{lem:upper-supply}

\section{83是一个统一够用的界，不是数学相变点}
三套计数都可以用 $p\ge83$ 一次保证，但其中两套在小素数上早已能用。偏移 $s\ge10$ 的情形，只需回到四个容量不等式逐项检查，小参数就能补齐。
\InsertProof{lem:upper-small}
这个引理的作用是缩小真正要查表的区域。不能因为某个统一证明选了83，就把所有小于83的参数都当作没有结构的特殊例子。

\section{直接造出n倍剪切：用一次守恒关系}
记来源为 $(i,j,k)$，新均值满足 $iA+jB+kC=pD$，全局还满足 $pA+pB+rC=0$。如果一个输出组的四项计数允许增加
\[\Delta_t=(ti-p, tj-p, tk-r, n-tp),
\]
那么项数不变，而输出总和增加
\[
t(iA+jB+kC)-(pA+pB+rC)+(n-tp)D=nD.
\]
所以新第一块均值增加 $nD/p=-nz'$，载体 $z'$ 不变。两条返回比较，恰好是 $(a',z')\mapsto(a'-nz',z')$，即 $U(-n)$。

这个根参数 $n$ 直接来自总和守恒。唯一还需检查的是：移动前后的两张输出表都不能超库存。五返回公式就是把这项检查与两个横向参数一起安排。
\InsertProof{lem:small-supply}

\section{剩余114组表，承担的是哪一个有限命题}
到此，无穷多个大 $p$ 已由公式覆盖，小 $p$ 的不少参数也由五返回公式覆盖。还保留114组物理根资源。它们各列一对或几对具体计数，检查非负、总数、迹零以及两个最大公因数。

查表以后，对任意有理参数的终止仍由同余输送证明。表里不枚举数值高度，不说“随机一千次都成功所以全部成功”。附录保留所有计数，你可以用上一讲三十四元的方式逐行读它。
